\chapter{Philotimic}

\section{Intro}
As I am moving on to college, one thing I note, beyond traditional internal self-authorship, there also exists the aspect of image. Image enforces behavior through wanting to be consistent, consistency builds virtue, image brings you together with like minded people, etc.

To continue to develop myself both internally and externally through consciously practiced drills and habitsis my goal. The highest human good in my world view. This area is going to be used both to explore the virtues and habits I wish to do and my way of practicing them in order to work on it. Then, log my activity.
\section{Overview}
\paragraph{Self-Perfection}

In this proccess, I will make mistakes. The purpose is to grow. I tend to forget about mistakes as soon as resolved, thus, a plan would be to write down daily errors, mistakes, way to improve, etc. I will likely put them in the drill log later in this chapter. Another option would be in the journal. I will figure out soon.

Some other drills would be to argue against my own position, Steelman other, and write about when my views change.

\paragraph{Outward image}
Before entering any professional room, take three seconds to check posture, breath, and warm neutral face.

Try to repeat people's names in order to help develop memory of them.

Whenever walking into a new room, come with a question or striking statement prepared.

I will also try to record 60 second introductions and work to improve them.

\paragraph{Speech, Vocabulary, Diction, and Rhetoric}

To improve my speech, I should work on trying to remove both filler words and vague quantifiers.

Three ways to practice this:

1: as I am developing my YouTube Videos, I will naturally develop these skills when working on the scripts reading them.

2: I can develop things that I wish to describe, explain, or define. Then, record myself speaking on them in order to see ways to naturally improve.

3: Lastly, take something that I said to another person and take some time each day to rewrite in a better form.

\paragraph{Academic Communication}

After stating a result, explain the scope and research.

when explaining or talking about a result. State the result, evidence for said statement, and finally end with limitations of my claim.

Write down paragraph summeries of research papers.

\paragraph{Assumptions and Cognative Habits}

Many people, including myself, rest ideas on assumptions we don't even know we have. I would like to work on this.

One way that I have already started is by writing my axiom book.

Another would be randomly taking random assumptions and either questioning them or proving them.

\paragraph{Bayes Reasoning}

Using Bayes reasoning to think more logically it Paramount. Getting to the habit is the problem.

One way would be to continuously be making predictions, guesses, conclusions, etc. Evaluate them using Bayes.

An example would be through every new room, evaluating people around me. Look at what they are doing, guess why, guess their job, age, height, and many more. When meeting people truly, try to come to a psychological profile of them of sorts.

\paragraph{Meta-analysis}

When evaluating my thinking, actions, plans, etc. Don't be vague, come up with an number.

\paragraph{Attention}

Evaluate wandering and conclude I am doing it. Don't feel shame and such, just fix and move on.

Try to be attentive even without deep work. Ponder, think on things, observe, to constantly be in a state of thought.

\paragraph{Observational awarness}

Many examples have already been said. Conclusions on people, observe things like exits, direction, fire extengushers.

\paragraph{Leadership}

The most obvious thing is to get into it and figure out ways to obtain leadership positions.

Some practice subjects could be mock meetings,

\paragraph{Dealines}

Every deadline or commitment, write it down to not forget. Get into the habit of using outlook and not relying on memory alone.

\paragraph{Body Language}

Straight posture, long gait, eyes straight, weight even, feet shoulder width apart, hands unclasped and seen, use hands rationally, and no stress responses.

Videos of me explaining things can be helpful in this regard.

\paragraph{Memory}

Practice memory skills. For instance, now I want to work on LaTeX snippets and Physics library, so I will work on that.

\paragraph{Desicion making}

Evaluate desicions. The information used to make them, my confidence of the results, actual results, how good/bad of a disusing it was, etc.

\paragraph{Teaching}

Practice teaching as much as possible; YouTube, personal videos, or writing.

\paragraph{Schedule}

I have now written schedules, I want to follow them and by my nature develop them for the future.

Once a week I will work on refining them.
\paragraph{Humor}

Avoid too much self-depricating humor as I have done in the past. Also, avoid making fun of others. Work on continuing my general ironic humor delivered in an ironical dry tone. I think that works well.

\paragraph{Writing}

Write more in this book. Even when I go to college, try to write something beyond lecture notes once a day.

I also plan to carry pocket note books with me.

\paragraph{Socratic Dialog}

In working on speaking, teaching, and debate. I would like to incorporate a kind of Socratic dialog with others. An obvious way to work on this is continuing to write essays. Another would be to actually practice it with others. How to slowly do that remains a challenge.

\section{Doing it}

\subsection{9/21/26}

\subsubsection{Morning Outlook}

As mentioned many times in previously, life is a project to develop ones self in virtue and understanding. I will be working on this project of mine in more detail today.

I plan to create a schedule based upon my earlier writing, I also wish to write a dialog on the ``paradox of altruism'' as I have been refering to it, essentially how can one be altruistic in a value system that puts altruism at its forefront.

I have been inspired by a program called the week of stoicism based upon early stoic teachings and thus am taking some ideas from them.
\subsubsection{Notes}
For my future actions, of the mentioned actions above I wish to prioritize

\begin{enumerate}
  \item Taking time to think without stimulation
  \item Socratic Dialog
  \item Beyond my traditional Socratic dialog where I wish to interrograte my thoughts not yet developed, I wish to interrogate my thoughts as they have been developed
  \item Evaluating ideas, actions, and mistakes at the end of the day
\end{enumerate}

I have written my Socratic dialog in a partially formulated way. Next I want to work on my dialectic section but I am not sure how to decide what to have a dialectic on.

An idea I wish to examine further is dual intentions, and how to apply that to myself and my actions.
\subsubsection{Review of the Day}
